\chapter{Project Context}
\label{chap:project-context}

This chapter frames the project before technical details begin. Present the host environment, the problem setting, the stakeholders, and the boundaries of the work.

\placeholderfigure{project-context-diagram.pdf}{Project context placeholder}{Insert a simple diagram showing the project environment, main actors, and expected outcome.}

\section{Internship Subject and Application Context}
Explain the subject in practical terms: the situation before the project, the need that motivated the work, and the application domain in which the solution will be used.

Focus on facts that help the reader understand the project. Avoid turning this section into a technical design description or a copy of the internship offer.

\section{Host Enterprise Presentation}
Present the host organization briefly: its domain, activities, relevant department or team, and why this environment is connected to the project.

Mention only information that supports the report context. Avoid promotional language, long corporate history, confidential details, or large catalogs of products and customers.

\section{Academic and Industrial Context}
Clarify the academic interest and practical value of the project. Connect the topic to engineering knowledge, research questions, professional practice, or operational constraints.

Avoid claiming scientific novelty without evidence. If the work is mainly applied engineering, say so clearly and explain its value.

\section{System Overview}
Give a high-level view of the existing system, process, or environment. Identify the main components, users, inputs, outputs, and dependencies.

\begin{table}[H]
    \centering
    \caption{Context summary template}
    \label{tab:context-summary-template}
    \small
    \begin{tabular}{L{0.28\textwidth}L{0.62\textwidth}}
        \toprule
        \textbf{Element} & \textbf{Guidance} \\
        \midrule
        Host environment & State where the work took place and which team or department was involved. \\
        \tabrowsep
        Problem setting & Describe the operational pain point or opportunity addressed by the project. \\
        \tabrowsep
        Main users & Identify who will use, review, maintain, or benefit from the result. \\
        \tabrowsep
        Constraints & Mention time, tools, data, hardware, access, standards, or organizational limits. \\
        \bottomrule
    \end{tabular}
\end{table}

\section{Stakeholder Needs}
List the needs of users, supervisors, maintainers, customers, reviewers, or other stakeholders. Explain conflicts or trade-offs when they influence later requirements.

Avoid inventing stakeholders only to fill space. Every need introduced here should connect to a requirement, design decision, or validation criterion later.

\section{Project Scope and Constraints}
Define what is included, what is excluded, and what assumptions limit the work. This section protects the report from unrealistic expectations.

Avoid hiding weak points. Clear boundaries make the validation chapter more credible.

\section{Conclusion}
Summarize the context in a few sentences and transition to the state of the art or literature review. Do not introduce new technical details in the conclusion.